\chapter{Objetivos}

\section{Objetivo general}

Dise\~nar y desarrollar un prototipo funcional de software que analice resonancias magn\'eticas lumbares sagitales para segmentar estructuras anat\'omicas relevantes, obtener un conjunto acotado de mediciones cuantitativas y generar una salida visual y estructurada editable que asista al profesional sin reemplazar su criterio cl\'inico.

\section{Objetivos espec\'ificos}

Para alcanzar el objetivo general, se definen los siguientes objetivos espec\'ificos:

\begin{enumerate}
    \item Relevar antecedentes acad\'emicos, datasets p\'ublicos, herramientas abiertas y soluciones comerciales relacionadas con el an\'alisis asistido de resonancias magn\'eticas lumbares.
    \item Identificar la brecha funcional que justifica el desarrollo de un prototipo acad\'emico orientado a segmentaci\'on, mediciones, visualizaci\'on y salida editable.
    \item Definir el alcance del MVP, especificando las estructuras anat\'omicas contempladas, las funcionalidades incluidas y las exclusiones cl\'inicas, t\'ecnicas y regulatorias.
    \item Analizar el flujo de trabajo y las necesidades de potenciales usuarios mediante actividades de user research orientadas a profesionales vinculados con resonancias magn\'eticas lumbares.
    \item Dise\~nar un pipeline de procesamiento que contemple carga, normalizaci\'on, preprocesamiento, segmentaci\'on, postprocesamiento, mediciones y visualizaci\'on de resultados.
    \item Implementar o adaptar un modelo de segmentaci\'on para identificar estructuras anat\'omicas en RM lumbar sagital, priorizando v\'ertebras, discos intervertebrales y canal espinal.
    \item Desarrollar un m\'odulo de mediciones cuantitativas simples derivadas de las m\'ascaras de segmentaci\'on generadas por el sistema.
    \item Dise\~nar una interfaz que permita visualizar la imagen original junto con las segmentaciones superpuestas y las mediciones asociadas.
    \item Generar una salida estructurada editable que organice los resultados del sistema sin constituir un informe cl\'inico autom\'atico.
    \item Evaluar t\'ecnicamente el desempe\~no del prototipo mediante m\'etricas de segmentaci\'on y comparaci\'on contra anotaciones de referencia del dataset utilizado.
    \item Complementar la validaci\'on t\'ecnica con una revisi\'on cualitativa profesional orientada a valorar plausibilidad anat\'omica, utilidad de las mediciones, claridad visual y legibilidad de la salida generada.
    \item Documentar las decisiones de dise\~no, limitaciones, resultados obtenidos y posibles l\'ineas de trabajo futuro.
\end{enumerate}
